\chapter*{Preface}

Volume~1 taught us to say exactly what is being proved and gave us polynomials, complex numbers and matrices. Volume~2 taught us to calculate with residues, primitive roots and Dirichlet characters. Those are not separate collections of tricks. For example, the character modulo $4$ sends the two invertible classes $[1]$ and $[3]$ to $1$ and $-1$, and it turns multiplication of classes into multiplication of signs. This volume gives a name to that behaviour: a \emph{homomorphism}. We will learn what such maps preserve and what they forget.

What does a square have in common with a shuffled deck of cards? You can change both without losing what makes them the same object. A rotation takes a square to itself; a shuffle takes a deck to another ordering of the same cards. We can perform two changes in succession, or undo a change. These simple moves lead to the language of groups.

Here is a harder question. The quadratic formula gives the roots of $ax^2+bx+c=0$ using arithmetic and a square root. There are longer formulas for cubics and quartics. Can we keep going? We will discover why the polynomial
\[
X^5-4X+2
\]
has no formula built from arithmetic and radicals. The obstruction is not that its roots are hard to approximate. It is that the permutations respecting their algebraic relations form $\mathfrak S_5$. To see why that matters, we first learn what a group can and cannot be broken into.

There is another puzzle. A cube has $24$ rotational symmetries, and
\[
24=1^2+1^2+2^2+3^2+3^2.
\]
The five squares are not a coincidence. They are the dimensions of the irreducible complex representations of the cube's rotation group, squared and added. By Chapter~11 we will be able to explain this identity, not just admire it.

The route starts with groups and their actions: why counting objects up to rotation needs more than division, why a normal subgroup is the right kind of subgroup to divide by, and what a permutation can tell us about an equation. Rings, modules and fields then provide the algebra we need to make Galois' idea precise. Finite fields give a different reward: we use polynomial arithmetic to recover a damaged message. Quaternions finally turn multiplication into three-dimensional rotation.

The appendices are optional detours, not missing prerequisites. One uses finite fields to construct Hadamard matrices. Another studies the positive eigenvector of a matrix. The other two show what happens when symmetry makes matrices easy to diagonalise or turns counting into linear independence. You can read each after its corresponding main chapter, or save them for later.

You will find small questions before some proofs, and graded exercises with hints and worked solutions after the chapters. Try an example before reading the abstract argument. When the notation gets dense, go back to the object it describes: a permutation, a coset, a polynomial, a rotation. The goal is to understand why the proof works, not to memorise its last line.

\emph{Order and Chance} follows the meeting of structure with uncertainty. Here we learn how much structure can be hidden in a change of viewpoint. Volume~4, \emph{Chance: The Logic of Uncertainty}, asks what can be learned when a change is chosen at random.
